\section{Naming Conventions}
\label{sec:naming-conventions}

To keep the report and diagrams consistent, we use the following naming conventions:

\begin{table}[!htbp]
  \centering
  \caption{Naming conventions and abbreviations.}
  \label{tab:naming-conventions}
  \begin{tabularx}{\linewidth}{@{}p{0.30\linewidth}X@{}}
    \toprule
    \textbf{Convention} & \textbf{Meaning} \\
    \midrule
    WezTerm & Always capitalize the product name as \textit{WezTerm}. \\
    Component names & Use title case for \textit{GUI Front End}, \textit{Multiplexer}, \textit{Mux Client}, \textit{Mux Server}, \textit{SSH Client}, and \textit{Terminal Emulation Core}. \\
    Mux & Use \textit{Mux} only in the names \textit{Mux Client} and \textit{Mux Server}; call the coordinating subsystem the \textit{Multiplexer}. \\
    Abbreviations & Capitalize \textit{GUI} (Graphical User Interface), \textit{SSH} (Secure Shell), \textit{PTY} (pseudo-terminal), and \textit{OS} (Operating System). \\
    External participants & In sequence diagrams, mark non-WezTerm participants with the UML stereotype \texttt{\guillemotleft{}external\guillemotright{}}, as for the Shell and SSH Server. In prose, identify them as external when relevant. \\
    Generic interface terms & Keep \textit{window}, \textit{tab}, \textit{pane}, \textit{domain}, and \textit{session} lowercase except at the start of a sentence. \\
    Dependency arrows & Solid arrows in the dependency diagram point from a component to a component it depends on; arrow labels state the dependency or flow. \\
    Sequence messages & Use concise, action-oriented lowercase labels, such as \textit{attach domain}, \textit{push changed lines}, and \textit{repaint changed lines}. Preserve capitalization for API or event identifiers such as \texttt{DetachDomain} and \texttt{PaneOutput}. \\
    \bottomrule
  \end{tabularx}
\end{table}
